\chapter{1979:Herbert C. Brown and Georg Wittig}

\label{chap:chemistry1979}

The Nobel Prize in Chemistry for the year 1979 was awarded jointly to Herbert C. Brown and Georg Wittig.

\textbf{Motivation:}

for their development of the use of boron- and phosphorus-containing compounds, respectively, into important reagents in organic synthesis

\section{Herbert C. Brown}

\subsection*{Early Life and Education}

Herbert C. Brown was born on 1912-05-22 in London, England.

\subsection*{Career and Major Research}

Brown immigrated to the United States with his family as a child and grew up in Chicago. He studied at Wright Junior College and the University of Chicago, where he received his doctorate. He taught at Wayne State University, now part of Wayne State University, and then at Purdue University, where he remained for the rest of his career. He discovered hydroboration, the addition of borane to alkenes, and developed the chemistry of organoboranes as versatile synthetic intermediates.

\subsection*{Nobel Prize-winning Work}

The work for which Herbert C. Brown was awarded the Nobel Prize in Chemistry in 1979 was described by the Nobel Committee as: "for their development of the use of boron- and phosphorus-containing compounds, respectively, into important reagents in organic synthesis".

Brown showed that hydroboration gives organoboranes that can be converted into a wide variety of functional groups with high selectivity and stereochemical control.

\subsection*{Significance and Impact}

Hydroboration and organoborane chemistry became among the most widely used methods in organic synthesis.

\subsection*{Later Career and Honors}

Brown remained at Purdue University until his retirement. He died on 2004-12-19 in West Lafayette, Indiana, USA.

\subsection*{Legacy}

Brown's boron chemistry is a standard toolkit of synthetic organic chemistry.

\subsection*{Modern Developments}

Brown's hydroboration reaction revolutionized organic synthesis by providing a mild, regio- and stereoselective way to convert alkenes into alcohols and other functional groups. Hydroboration and the related organoborane chemistry are now standard tools in the synthesis of pharmaceuticals and natural products, and boron reagents feature prominently in modern cross-coupling and borylation chemistry.

\subsection*{Selected Publications}

"Hydroboration" (1962)

\clearpage

\section{Georg Wittig}

\subsection*{Early Life and Education}

Georg Wittig was born on 1897-06-16 in Berlin, Germany.

\subsection*{Career and Major Research}

Wittig studied at the University of Marburg, where he received his doctorate. He held professorships at Marburg, Braunschweig, Tübingen, and Heidelberg. He discovered the rearrangement of ylides and the reaction that bears his name, in which a phosphorus ylide reacts with a carbonyl compound to form an alkene.

\subsection*{Nobel Prize-winning Work}

The work for which Georg Wittig was awarded the Nobel Prize in Chemistry in 1979 was described by the Nobel Committee as: "for their development of the use of boron- and phosphorus-containing compounds, respectively, into important reagents in organic synthesis".

The Wittig reaction allows chemists to place a carbon-carbon double bond at a defined position in a molecule, with predictable stereochemistry.

\subsection*{Significance and Impact}

The Wittig reaction became one of the most valuable methods for synthesizing alkenes, including the natural product vitamin A.

\subsection*{Later Career and Honors}

Wittig remained at the University of Heidelberg until his retirement. He died on 1987-08-26 in Heidelberg, Germany.

\subsection*{Legacy}

The Wittig reaction is a fundamental tool for constructing carbon-carbon double bonds in organic synthesis.

\subsection*{Modern Developments}

Wittig's discovery of the olefination reaction that bears his name provided a reliable method for constructing carbon-carbon double bonds, a transformation now used routinely in the synthesis of pharmaceuticals, natural products, and materials. The Wittig reaction remains a cornerstone of organic synthesis and has been extended into modern variants such as the Horner–Wadsworth–Emmons reaction.

\subsection*{Selected Publications}

"Zur Reaktionsweise des Pentaphenyl-phosphors und einiger Derivate" (Liebigs Annalen der Chemie, 1953)

\clearpage

\section*{Chapter Summary}

The 1979 prize recognized Brown and Wittig for developing boron- and phosphorus-containing reagents into important tools of organic synthesis. Their methods are essential to modern synthetic chemistry.

